\noindent\textbf{Text-attributed graph learning.}
TAGs encode node semantics through text and relations through edges. Existing
benchmarks compare language models, graph neural networks, and their joint
training~\cite{tag_benchmark}. Common graph encoders include convolution,
attention, message passing, and inductive neighborhood aggregation
~\cite{gcn,gat,mpnn,graphsage,gin}. Earlier network representation methods
use random-walk neighborhoods to obtain compact node features
~\cite{deepwalk,node2vec}, while text encoders and multimodal
alignment models provide semantic representations~\cite{bert,sbert,clip}.
GraphCLIP learns transferable representations by contrasting graphs with text
summaries~\cite{graphclip}, and graph self-supervision provides related ways
to improve representation quality~\cite{dgi,graphcl}. Continual adaptation
coordinates modality updates to mitigate forgetting~\cite{g2lora}. These
studies focus on learning and adapting models. We study the training data
retained between updates, constructing compact memories that contain both
numerical features and real text for subsequent encoding.

\noindent\textbf{Graph condensation.}
Dataset condensation has used gradient matching, distribution matching, and
training-trajectory matching to preserve the behavior of a larger training
set~\cite{dataset_gradient,dataset_distribution,trajectory_matching}.
Graph condensation extends these ideas to relational data~\cite{gcond,sfgc}
and has developed benchmarks, inductive node mappings, and open-world
settings~\cite{gcondenser,mcond,opengc}. Structure can guide feature
construction even when the resulting training set stores no synthetic edges.
Our setting additionally requires readable text for every memory item. The
learned support matrix coordinates the contribution of current nodes to
numerical features and the selection of a real text source, giving these two
components a common construction variable.

\noindent\textbf{TAG and multimodal graph condensation.}
TAGSAM combines subgraph text selection with attribute matching~\cite{tagsam}. STAD combines cluster compression with text generation~\cite{stad}, and TaLK optimizes continuous token embeddings~\cite{talk}. DANCE uses training feedback for dynamic neighborhood and text selection in federated learning~\cite{dance}. SR-GM addresses modality conflicts through gradient decoupling and structural regularization~\cite{srgm}. These methods establish several ways to couple text and graph compression. Our focus is constructing successive memory increments after earlier raw data become unavailable. Within each increment, real text selection receives surrogate gradients through the support used for feature aggregation, while historical memory supplies context for the construction objective.

\noindent\textbf{Continual graph learning and evolving condensation.}
Continual learning studies how to preserve earlier knowledge while new data
arrive, using replay, parameter regularization, or distillation
~\cite{icarl,ewc,lwf,continual_survey}. Dynamic graph models instead update
representations as nodes, attributes, and edges evolve~\cite{evolvegcn,tgn,
dynamic_survey}. Class-incremental graph learning also exposes the tension
between retaining earlier classes and adapting to new ones~\cite{qi2024}.
In graph condensation, replay methods address historical
coverage and class representativeness~\cite{dslr}, CaT reduces the imbalance
between incoming data and replay through condensation and training in memory
~\cite{cat}, PUMA incorporates unlabeled information and reduces repeated
propagation and condensation costs~\cite{puma}, and GECC inherits clustering
states of aggregated features for evolving condensation~\cite{gecc}. These
approaches motivate a compact state that can persist across updates. We extend
that interface to paired numerical features and readable text, constructing
new items from the current increment and historical memory. This addresses
the additional information that a graph and text learner requires after the
original increment is released.
